\documentclass[11pt,a4paper]{article}
\usepackage{../shared/hanzocover}
\usepackage[utf8]{inputenc}
\usepackage[T1]{fontenc}
% Shared preamble for the compute-market series.
% One style, one vocabulary, inputted by every paper so the series reads as one voice.
\usepackage[margin=1in]{geometry}
\usepackage{booktabs}
\usepackage{amsmath}
\usepackage{amssymb}
\usepackage{amsthm}
\usepackage{graphicx}
\usepackage{xcolor}
\usepackage{hyperref}
\hypersetup{colorlinks=true, linkcolor=blue, urlcolor=blue, citecolor=blue}
\usepackage{enumitem}
\usepackage{listings}
\usepackage{array}
\usepackage{tabularx}
\usepackage{siunitx}
\sisetup{group-separator={,}, group-minimum-digits=4}

\definecolor{kw}{rgb}{0.13,0.29,0.53}
\definecolor{cmt}{rgb}{0.40,0.40,0.40}
\definecolor{str}{rgb}{0.16,0.50,0.16}
\lstset{
  basicstyle=\ttfamily\footnotesize,
  keywordstyle=\color{kw}\bfseries,
  commentstyle=\color{cmt}\itshape,
  stringstyle=\color{str},
  showstringspaces=false,
  breaklines=true,
  columns=fullflexible,
  frame=single,
  framesep=4pt,
}

\newtheorem{definition}{Definition}
\newtheorem{proposition}{Proposition}
\newtheorem{remark}{Remark}

\newcommand{\code}[1]{\texttt{#1}}
\newcommand{\repo}[1]{\texttt{#1}}

% Series vocabulary — used identically in every paper.
\newcommand{\job}{\ensuremath{J}}
\newcommand{\bundle}{\ensuremath{R}}
\newcommand{\cost}{\ensuremath{\mathcal{C}}}

\tolerance=800
\emergencystretch=3em


\title{Compute With Memory:\\Adapter Residency as a Price Term}

\author{
Zach Kelling\thanks{Corresponding author: research@hanzo.ai} \\
\textit{Hanzo Industries Inc (Techstars '17)} \\
Los Angeles, California \\
\texttt{https://hanzo.ai}
}

\date{2026}

\begin{document}
\hanzocoverpage

\begin{abstract}
Proof of work rests on an assumption that is easy to miss because it is never
stated: hashpower is memoryless. Any miner is interchangeable with any other
miner at the same rate, and no miner is cheaper for one particular block than
for another. That assumption is what permits a single scalar difficulty.

Personalized agents break it. When an agent carries its own low-rank weight
delta, a provider that has served that agent before is permanently cheaper for
that agent than an otherwise identical provider that has not. Compute acquires
memory, the supply side becomes non-fungible per buyer, and cost stops being a
function of the endpoint and becomes a function of the path taken through
resource space. That is precisely the condition under which a scalar price
fails and a Hamiltonian formulation earns its keep.

We quantify the effect. A rank-16 adapter on a 70B model is
\SI{168}{\mega\byte} against a \SI{140}{\giga\byte} base --- a residency
asymmetry of roughly three orders of magnitude. We derive the resulting cost
functional, show that the na\"ive one-process-per-adapter deployment costs
\num{175} accelerators where a batched deployment costs two, and state the
single serving primitive that separates those two regimes. We report the exact
state of that primitive in the Hanzo inference engine, including what is absent.
\end{abstract}

\tableofcontents
\newpage

\section{The memorylessness assumption}

Bitcoin's difficulty adjustment is a scalar because the resource it prices is
one-dimensional and history-free. Write $h_i$ for miner $i$'s hash rate. The
expected time for miner $i$ to find a block at difficulty $D$ is $D/h_i$,
independent of which block, which transactions it contains, and everything
miner $i$ has ever done before. Two consequences follow, and both are load-bearing:

\begin{enumerate}[label=(\alph*)]
\item \emph{Fungibility.} Miners at equal $h$ are perfect substitutes. There is
no buyer-specific cost.
\item \emph{Statelessness.} The cost of the next unit of work does not depend on
the sequence of previous units. There is no path dependence.
\end{enumerate}

Together, (a) and (b) collapse the market to a single price per unit hash and a
single difficulty. Neither survives contact with personalized inference.

\section{Where the memory comes from}

\subsection{Low-rank deltas}

A low-rank adaptation replaces a frozen projection $W \in \mathbb{R}^{d_{\mathrm{out}} \times d_{\mathrm{in}}}$
with
\begin{equation}
W' x = W x + s \cdot B A x,
\qquad
A \in \mathbb{R}^{r \times d_{\mathrm{in}}},\;
B \in \mathbb{R}^{d_{\mathrm{out}} \times r},\;
s = \frac{\alpha}{r},
\end{equation}
where $r \ll \min(d_{\mathrm{in}}, d_{\mathrm{out}})$. The pair $(A,B)$ is the
agent's private state. Everything else is shared.

This is the shape implemented in \repo{hanzo-train/src/lora.rs}, where the base
is a frozen \code{Linear} and the delta is a separate tracked
\code{LoraDelta \{ name, a, b, scale \}} carrying $A$ as $(r, d_{\mathrm{in}})$
and $B$ as $(d_{\mathrm{out}}, r)$.

\subsection{Parameter count and size}

Adapting $T$ projections in each of $L$ layers of a model with hidden dimension
$d$ costs
\begin{equation}
|\theta_{\text{adapter}}| \;=\; 2\,r\,L\,T\,d
\end{equation}
parameters when the adapted projections are square in $d$. At \code{bf16} the
byte cost is twice that. Table~\ref{tab:sizes} gives the values that matter.

\begin{table}[h]
\centering
\begin{tabular}{llrrr}
\toprule
Model class & Targets & $r$ & Parameters & Size (\code{bf16}) \\
\midrule
7B ($d{=}4096$, $L{=}32$)  & q,k,v,o        & 8  & \num{8388608}   & \SI{16.8}{\mega\byte} \\
7B                         & q,k,v,o        & 16 & \num{16777216}  & \SI{33.6}{\mega\byte} \\
7B                         & q,k,v,o + mlp  & 16 & \num{29360128}  & \SI{58.7}{\mega\byte} \\
7B                         & q,k,v,o        & 64 & \num{67108864}  & \SI{134.2}{\mega\byte} \\
70B ($d{=}8192$, $L{=}80$) & q,k,v,o        & 8  & \num{41943040}  & \SI{83.9}{\mega\byte} \\
70B                        & q,k,v,o        & 16 & \num{83886080}  & \SI{167.8}{\mega\byte} \\
70B                        & q,k,v,o        & 64 & \num{335544320} & \SI{671.1}{\mega\byte} \\
\bottomrule
\end{tabular}
\caption{Adapter parameter counts and sizes.}
\label{tab:sizes}
\end{table}

\subsection{The asymmetry}

Set against those figures, the base model is \SI{140}{\giga\byte} at
\code{bf16} for a 70B, or roughly \SI{40}{\giga\byte} quantized to four bits.
The ratio that governs everything downstream is
\begin{equation}
\frac{|\text{base}|}{|\text{adapter}|} \;\approx\; 10^{3}.
\end{equation}

An adapter is nearly free to move; the base it attaches to is, for practical
purposes, immovable within any interactive latency budget. At
\SI{10}{\giga\bit\per\second} a \SI{140}{\giga\byte} transfer takes
\SI{112}{\second}. Against a \SI{3}{\second} service objective that is not
expensive --- it is infeasible. Feasibility and price are therefore entangled
through residency, which is the technical heart of this paper.

\section{The cost functional}

\begin{definition}[Residency state]
For a provider $\bundle$ at time $t$, let $\beta(\bundle,t) \in \{0,1\}$
indicate whether the base model required by job $\job$ is resident in
accelerator memory, and $\alpha(\bundle,\job,t) \in \{0,1\}$ whether $\job$'s
adapter is resident.
\end{definition}

The allocation objective in the heterogeneous market is then

\begin{equation}
\cost(\job,\bundle) \;=\;
\underbrace{p(\bundle)\,\tau(\job,\bundle)}_{\text{compute}}
\;+\; \lambda_{L}\,\ell
\;+\; \lambda_{F}\,\phi
\;+\; \lambda_{V}\,v
\;+\; \underbrace{\mu_{1}\big(1-\beta\big) \;+\; \mu_{2}\big(1-\alpha\big)}_{\text{cold-start}},
\label{eq:objective}
\end{equation}

with $\ell$ latency, $\phi$ failure risk, $v$ verification cost, and
\begin{equation}
\frac{\mu_{1}}{\mu_{2}} \;\approx\; \frac{|\text{base}|}{|\text{adapter}|} \;\approx\; 10^{3}.
\end{equation}

The last two terms are what a stateless spot market omits, and they are not a
correction --- at interactive deadlines $\mu_1$ is frequently large enough to
make the constraint binding rather than the price.

\begin{proposition}[Cheapest capable is wrong]
Let $\bundle_{1}$ be an accelerator priced at $p_{1}$ with $\beta = 1$, and
$\bundle_{2}$ priced at $p_{2} \ll p_{1}$ with $\beta = 0$. If the deadline
$\Delta$ satisfies $\Delta < |\text{base}| / \text{bw}(\bundle_{2})$, then
$\bundle_{2}$ is infeasible regardless of $p_{2}$, and the correct allocation is
$\bundle_{1}$ at arbitrarily large price ratio.
\end{proposition}

The worked instance: an accelerator at \$4/hour holding the model beats one at
\$0.50/hour that must first fetch \SI{140}{\giga\byte}. Under a
\SI{3}{\second} objective the cheaper machine is not in the feasible set at all.

\section{Affinity, and what it implies for the market}

Because $\mu_{1} \gg \mu_{2}$, agents develop \emph{affinity} for the nodes that
already hold their base model, and their adapter follows cheaply. This produces
structure a spot market cannot express:

\paragraph{Stickiness is efficient, not rent-seeking.} Routing an agent back to
the same provider is the cost-minimising choice, not a lock-in artifact. A
market design that forces uniform random allocation across capable providers
destroys real surplus.

\paragraph{Eviction is a pricing decision.} An accelerator holding $k$ resident
adapters faces a knapsack problem whenever a new agent arrives. The shadow price
of a residency slot is the expected discounted future value of the agent
occupying it, which is exactly the object a continuous market can quote and a
discrete order book cannot.

\paragraph{Residency is self-proving.} A provider cannot profitably lie about
holding an agent's weights, because the lie is revealed by the deadline it then
misses. This is unusual and useful: the market needs no attestation mechanism
for this particular claim.

\paragraph{State rent is a new instrument.} An agent may pay a small continuous
fee to keep its adapter resident somewhere, converting a stochastic cold-start
cost into a predictable stream. This is an option on locality. It has no
analogue in a stateless compute market and it is a revenue line that does not
exist today.

\section{Capacity: how many agents fit}

The reason this matters economically is that adapters are small enough that
personalization is not, in fact, expensive --- once the delta stops dragging a
base model behind it.

\begin{table}[h]
\centering
\begin{tabular}{llrr}
\toprule
Base & Free memory & Adapter & Resident agents \\
\midrule
70B \code{bf16} on \SI{180}{\giga\byte} & \SI{40}{\giga\byte} & \SI{17.2}{\mega\byte}  & \num{2326} \\
70B \code{bf16} on \SI{180}{\giga\byte} & \SI{40}{\giga\byte} & \SI{137.4}{\mega\byte} & \num{291} \\
70B 4-bit on \SI{80}{\giga\byte}        & \SI{40}{\giga\byte} & \SI{17.2}{\mega\byte}  & \num{2326} \\
7B \code{bf16} on \SI{80}{\giga\byte}   & \SI{65}{\giga\byte} & \SI{17.2}{\mega\byte}  & \num{3779} \\
\bottomrule
\end{tabular}
\caption{Resident agents beside one shared base, ignoring KV cache.}
\label{tab:capacity}
\end{table}

Against that, the cost of \emph{not} sharing the base:

\begin{table}[h]
\centering
\begin{tabular}{lrr}
\toprule
100 personalized agents, 70B \code{bf16} & Memory & Accelerators \\
\midrule
One process per adapter & \SI{14000}{\giga\byte} & \num{175} \\
Batched multi-adapter   & \SI{141.7}{\giga\byte} & \num{2} \\
\bottomrule
\end{tabular}
\caption{The duplication cost, at \SI{80}{\giga\byte} per accelerator.}
\label{tab:duplication}
\end{table}

An \(\approx 87\times\) reduction in the capital required to serve a
personalized population is not an optimization. It sets the cost floor for the
entire category, and whoever ships it first defines that floor for everyone else.

\section{What exists, and the one thing that does not}

We state the current position precisely, because the gap is narrow and specific
and it is worth no one's time to mistake it for a research problem.

\subsection{Training: complete}

\repo{hanzo-train} is \num{2215} lines implementing trainable LoRA over frozen
bases, AdamW, distillation, a data loader, and PEFT-format serialization whose
tensor names match the inference-side merge. It is exposed over the serving API:

\begin{lstlisting}
POST   /v1/training/clients                       base model + LoRA injection
GET    /v1/training/clients[/{id}]                list / inspect (+ loss history)
POST   /v1/training/clients/{id}/forward_backward accumulate gradients
POST   /v1/training/clients/{id}/optim_step       apply AdamW
POST   /v1/training/clients/{id}/sample           decode from current base+delta
POST   /v1/training/clients/{id}/save_weights     write the adapter for inference
DELETE /v1/training/clients/{id}                  drop client, free memory
\end{lstlisting}

An agent can be given its own weights today, end to end.

\subsection{Serving: co-residency yes, selection no}

The inference side carries \repo{hanzo-engine/src/lora/loralinear.rs} and
\code{qloralinear.rs}, an X-LoRA implementation under
\repo{hanzo-engine/src/xlora\_models/}, and AnyMoE. X-LoRA is significant here
for a reason unrelated to its purpose: it loads many adapters into one process
simultaneously. \emph{Co-residency is solved.} What X-LoRA does with those
adapters --- mix all of them for every token through a learned classifier --- is
not what a market needs, but the memory management, shape handling and loading
are done.

What is absent is selection:

\begin{itemize}[leftmargin=1.4em]
\item \code{preload\_adapters} is a load-time \code{HashMap}, threaded through
      the loaders as \code{Option<HashMap<String, (ShardedVarBuilder, LoraConfig)>>}.
\item \code{with\_lora(adapters: Vec<String>)} is a pipeline-construction argument.
\item \code{hanzo\_quant::lora::merge\_lora\_weights} folds the delta \emph{into}
      the base weight at load, which destroys per-adapter separability by design.
\item The OpenAI-compatible request in \repo{hanzo-server-core/src/openai.rs}
      carries no adapter field.
\end{itemize}

One process therefore serves one adapter configuration, which is exactly the
regime that costs \num{175} accelerators in Table~\ref{tab:duplication}.

\subsection{The missing primitive}

Batched multi-adapter serving requires that a single forward pass over a batch
of sequences apply a \emph{different} delta per sequence:
\begin{equation}
y_{i} \;=\; W x_{i} \;+\; s_{a(i)} \, B_{a(i)} A_{a(i)} x_{i},
\qquad a(i) \in \{1,\dots,k\},
\end{equation}
where $a(i)$ is the adapter index of sequence $i$ in the batch. Implemented as a
segmented gather-matmul, the added cost over the shared base is
$O(r(d_{\mathrm{in}}+d_{\mathrm{out}}))$ per token --- negligible against the
$O(d_{\mathrm{in}}d_{\mathrm{out}})$ base projection at $r \le 64$.

The unmerged forward already exists in the tree, on the training side:
\code{LoraLinear} in \repo{hanzo-train/src/lora.rs} keeps \code{base} frozen and
\code{delta: Option<LoraDelta>} separate, computing $s\,(x A^{\!\top}) B^{\!\top}$
as its own path with tensor layouts documented as matching the inference merge.
The work is to make that path batch-aware over $a(i)$, add an adapter field to
the request, and place a cache with eviction in front of it. Every piece has a
reference implementation in the same repository.

\section{The same trick one layer down}

The pattern in this paper --- a large frozen shared object plus a tiny per-tenant
delta --- is not confined to model weights. Our pricing layer uses it too, and the
symmetry is worth naming because it suggests the shape is general rather than an
artifact of how transformers happen to be built.

There, tenants differ in how they trade price against latency and quality. The
straightforward implementation is a full-precision preference table per tenant,
which does not survive scale. Instead the tenant's deviation is carried as a
one-bit sign delta over a shared base quote table --- one bit per regime and
resource cell, with per-layer scale factors and higher-precision bias terms
retained. The compression against a single-precision table is thirty-two-fold,
which is the difference between holding thousands of tenant tables in memory and
holding millions.

Two layers, same structure: the shared object is large and effectively immovable,
the per-tenant object is tiny, and the economics of the whole system fall out of
that ratio. Where the analogy stops is instructive. At the pricing layer the
shared table is small enough that duplicating it per tenant is merely wasteful; at
the model layer the shared object is a hundred and forty gigabytes and duplicating
it is the difference between two accelerators and a hundred and seventy-five. The
asymmetry is what makes selection --- rather than duplication --- the load-bearing
primitive at the model layer specifically.

\section{Consequence for a Hamiltonian market}

With stateless homogeneous inference, a spot market for accelerator-hours is
close to sufficient; the resource manifold is nearly a line and a price vector
nearly spans it. Personalization is what makes the manifold genuinely coupled:
the cost of serving a request depends on the trajectory the system has already
taken, because that trajectory determines what is resident where.

This is the concrete content of treating compute as a coupled resource system
rather than a family of independent order books. The coupling is not between
device types --- those really are close to substitutes at equal capability. The
coupling is between \emph{a buyer's history} and \emph{a provider's current
state}, and it is strong, asymmetric, and continuously varying. A scalar
difficulty cannot express it. A Hamiltonian in $(q,p)$ over occupancy and price
can.

\section*{Availability}

The inference engine, training crate, router and model catalog discussed here
are open source under permissive licence. The adapter arithmetic in
Tables~\ref{tab:sizes}--\ref{tab:duplication} is reproducible from the formulas
given; no proprietary measurement is involved.

\end{document}
